\section*{Chapter \ref{ch:g2:where-animals-live} --- Where Animals Live}

\begin{solution}{exo:g2:where-animals-live:1}
The kind of place where an \omterm{def:g1:animals-around-us:animal}{animal} naturally lives. It finds there
its food, water, shelter, and a safe spot for its young.
\end{solution}

\begin{solution}{exo:g2:where-animals-live:2}
Pond: frog, dragonfly (or stickleback, pond snail). Forest:
squirrel, woodpecker (or deer, wood ant). Meadow: rabbit,
grasshopper (or butterfly, field mouse).
\end{solution}

\begin{solution}{exo:g2:where-animals-live:3}
Stickleback: the pond. Woodpecker: the forest. Grasshopper: the
meadow. Woodlouse: under the stones of an old wall (damp hidden
corners).
\end{solution}

\begin{solution}{exo:g2:where-animals-live:4}
The frog: webbed feet that push water. The squirrel: gripping
claws and a big balancing tail.
\end{solution}

\begin{solution}{exo:g2:where-animals-live:5}
Wide digging shovel-hands; \omterm{def:g1:animals-around-us:covering}{fur} lying flat both ways so soil never
ruffles it; \omterm{def:g1:the-five-senses:sense}{eyes} it hardly needs in the dark.
\end{solution}

\begin{solution}{exo:g2:where-animals-live:6}
Woodpecker: high on the trunks. Squirrel: in the \omterm{def:g1:plants-around-us:tree}{branches}. Deer:
among the bushes. Wood ant: on the ground, in its heaped nest.
\end{solution}

\begin{solution}{exo:g2:where-animals-live:7}
They disappear with the pond --- without the \omterm{def:g2:where-animals-live:habitat}{habitat} they lose
food, shelter and breeding place all at once. So protecting
\omterm{def:g1:animals-around-us:animal}{animals} starts with protecting the places they live.
\end{solution}

\begin{solution}{exo:g2:where-animals-live:8}
Open answer. Even a single turned stone often shelters woodlice,
ants, a beetle, sometimes a centipede --- small \omterm{def:g2:where-animals-live:habitat}{habitats} are
surprisingly full.
\end{solution}

\begin{solution}{exo:g2:where-animals-live:9}
Webbed feet (paddling), water-shedding \omterm{def:g1:animals-around-us:covering}{feathers}, a sifting beak:
all three point to water --- a pond or lake \omterm{def:g2:where-animals-live:habitat}{habitat}. The \omterm{def:g1:animals-around-us:animal}{animal}
is very likely a duck.
\end{solution}

\begin{solution}{exo:g2:where-animals-live:10}
The frog's equipment --- webbed feet, moist \omterm{def:g1:the-five-senses:sense}{skin}, a body made for
water and damp banks --- is useless in a dry box, and its \omterm{def:g2:where-animals-live:habitat}{habitat}
gave it more than food: water for its \omterm{def:g1:the-five-senses:sense}{skin} and its \omterm{def:g2:animals-grow-up:born}{eggs}, hiding
places, other frogs. Kindness to a wild \omterm{def:g1:animals-around-us:animal}{animal} is leaving it in
--- or returning it to --- its \omterm{def:g2:where-animals-live:habitat}{habitat}.
\end{solution}
